
% =====================================================================
\section{Robustness}
\label{sec:robust}
% =====================================================================

\subsection{The placebo}

\begin{figure}[htbp]
\centering
\includegraphics[width=0.85\linewidth]{f_placebo.pdf}
\caption{The effective-spread regression re-run with the measure rebuilt on each
last price digit in turn. The round-price digit is highlighted; bars are 95\%
confidence intervals.}
\label{fig:placebo}
\end{figure}

\begin{table}[htbp]
\centering
\caption{Placebo test: the same regression run on every last price digit}
\label{tab:placebo}
\small
\begin{threeparttable}
\begin{tabular}{@{}lrrr@{}}
\toprule
Last digit & Coefficient & SE & $t$ \\
\midrule
0 & 0.0192*** & (0.0075) & 2.58 \\
1 & -0.0128 & (0.0092) & -1.40 \\
2 & 0.0057 & (0.0103) & 0.55 \\
3 & 0.0043 & (0.0085) & 0.50 \\
4 & -0.0023 & (0.0099) & -0.24 \\
5 & -0.0017 & (0.0092) & -0.19 \\
6 & -0.0008 & (0.0097) & -0.09 \\
7 & -0.0088 & (0.0094) & -0.94 \\
8 & -0.0120 & (0.0091) & -1.32 \\
9 & -0.0089 & (0.0098) & -0.91 \\
\bottomrule
\end{tabular}
\begin{tablenotes}[flushleft]\footnotesize
\item Outcome is the log effective spread in the dynamic specification of Table~\ref{tab:rq2}, with the round-price share replaced by the share of buy-initiated large-trade volume at each last digit in turn. If the round-price result reflects stale limit orders at psychologically salient prices, digits other than zero should carry no comparable information. Digit five is marked as semi-focal because half-way points attract some of the same behaviour.
\end{tablenotes}
\end{threeparttable}
\end{table}


This is the test that could have ended the study. The measure is a share of
volume at a particular last digit. If a share of volume at \emph{any} last
digit predicts liquidity equally well, then nothing here is about round numbers
--- it is an artefact of how a concentration measure behaves, and the whole
interpretation collapses.

Rebuilding the measure on each of digits one through nine and re-running the same
regression, the round-price coefficient is nan
($t = nan$), and 0 of the nine placebo digits reach
conventional significance. Digit five is flagged separately in
Table~\ref{tab:placebo} because half-way points are semi-focal and attract some
of the same behaviour --- treating it as a pure placebo would be stacking the
deck.

\subsection{Subsamples and an alternative estimator}

\begin{table}[htbp]
\centering
\caption{Robustness of the clustering--liquidity relationship}
\label{tab:robust}
\small
\begin{threeparttable}
\begin{tabular}{@{}lrrr@{}}
\toprule
Specification & Coefficient & SE & Stock-days \\
\midrule
Baseline (dynamic) & 0.0188** & (0.0075) & 74,034 \\
Size Q1 & 0.0140 & (0.0156) & 15,062 \\
Size Q2 & 0.0313** & (0.0127) & 14,999 \\
Size Q3 & 0.0217 & (0.0146) & 15,099 \\
Size Q4 & -0.0119 & (0.0174) & 14,743 \\
Size Q5 & 0.0207 & (0.0241) & 14,054 \\
Period H1 (2025-01 to 2025-03) & 0.0164* & (0.0096) & 36,842 \\
Period H2 (2025-03 to 2025-04) & 0.0037 & (0.0103) & 37,192 \\
Excluding 2025-04-03 to 2025-04-11 & 0.0189** & (0.0079) & 67,970 \\
Five lags of the outcome & 0.0124* & (0.0073) & 70,014 \\
0.1-yen-tick days only & -0.0033 & (0.0613) & 2,057 \\
\bottomrule
\end{tabular}
\begin{tablenotes}[flushleft]\footnotesize
\item Each row re-estimates the dynamic specification of Table~\ref{tab:rq2} with the log effective spread as the outcome, on the stated subsample or with the stated change. Size quintiles are formed once per stock on median market capitalisation. The period rows split the panel at its median trading date, derived from the sample rather than from a calendar. The event-window row removes the named week. The five-lag row replaces the single lag of the outcome with five, so the coefficient reads against a richer version of the outcome's own history. $^{*}p<0.1$, $^{**}p<0.05$, $^{***}p<0.01$.
\end{tablenotes}
\end{threeparttable}
\end{table}


Table~\ref{tab:robust} re-estimates the headline specification on subsamples.
The size-quintile rows matter for interpretation: the mechanism is about
retail-heavy, less closely watched stocks, so the relationship should be stronger
in smaller ones. The half-year rows test whether a single episode is doing the
work. The 0.1-yen row estimates on the days excluded from the main sample.

A Fama--MacBeth cross-check estimates the relationship separately in each daily
cross-section and averages, with Newey--West standard errors over the resulting
series. It gives -- with $t = --$ over
0 cross-sections. It is reported because it makes completely
different assumptions about the correlation structure than two-way clustering
does, and agreement between them is worth more than either alone.

\paragraph{Verdict.}
The result is specific to the round-price digit rather than to the shape of the
measure, and it survives changes in subsample and estimator.
